\chapter{Discussion}
\label{ch:discussion}

This chapter brings the two parts together. \Cref{sec:synthesis}
translates the design space of \Cref{ch:design} and the measurements of
\Cref{ch:results} into a recommended configuration and shows how the
recommendation changes with venue type. \Cref{sec:testable2} states
precisely which theoretical predictions this design tested, which it
confirmed by construction, and which it could not reach.
\Cref{sec:positioning} places the findings against the existing
literature, and \Cref{sec:validity} sets out the threats to validity.

\section{From Design Space to Recommendation}
\label{sec:synthesis}

\subsection{The Structure of the Problem}

The recurring finding of \Cref{ch:design} is that the FBA design space
contains almost no free improvements. Each block offers options that
shift competition from one variable to another rather than removing it.
Randomising the close removes the timing target and adds planning cost.
Opacity removes intent exposure and adds bid shading. Pro rata removes
the timing channel from rationing and adds a size channel. Segregation
lets makers quote tighter against benign flow and splits the pool.

The one genuine structural gain is the pair of rules that constitutes
the mechanism itself: discrete accumulation plus uniform pricing. That
pair removes the value of a marginal speed advantage without requiring
anything from participants and without depending on how any other block
is configured. Everything else in the design space is a choice about
which residual distortion to accept.

This has a practical consequence for how an exchange operator should
approach the problem. The blocks should not be optimised independently,
because they interact. \Cref{sec:disclosure} showed the sharpest
instance: an open book during accumulation strengthens the incentive to
submit late, which raises the value of a randomised close, whereas a
sealed-bid design removes most of the informational motive for waiting
and makes a fixed boundary defensible. Disclosure and interval design
must be chosen jointly. The same is true of interval length and
instrument liquidity, since a short interval in a thin instrument
produces empty batches and degenerates into a delayed continuous
market.

\subsection{A Recommended Baseline}

\Cref{tab:recommendation} states a baseline configuration for a liquid
instrument on a centralised venue, with the reasoning behind each
choice. The recommendation follows from the theory of \Cref{ch:design}
and is to be confirmed or revised against the results of
\Cref{ch:results} once they are populated; the final column records
which measurement bears on the choice.

\begin{table}[htbp]
\centering
\caption[Recommended baseline configuration]{A recommended baseline FBA
configuration for a liquid instrument on a centralised venue.}
\label{tab:recommendation}
\small
\begin{tabularx}{\textwidth}{@{}p{2.3cm} p{3.0cm} L p{2.4cm}@{}}
\toprule
\textbf{Block} & \textbf{Recommendation} & \textbf{Reasoning} &
\textbf{Evidence}\\
\midrule
Interval & 50--250\,ms, with a short randomised extension &
Long enough that the ratio of latency dispersion to $\tau$ is
negligible; short enough that delay is immaterial to any participant
other than a market maker managing inventory & \Cref{tab:sweep};
empty-batch share\\
Disclosure & Indicative clearing price and aggregate imbalance; no
order-level book & Retains the coordination benefit of transparency
without exposing individual orders; the regime equity auctions already
use at scale & Not testable here (P7)\\
Segregation & Unified pool by default; intent segregation only where
maker flow is demonstrably being adversely selected & Segregation splits
liquidity, which is a certain cost against an uncertain benefit; the
DFBA argument is strongest where the unified pool shows persistent
maker-versus-maker crossing & \Cref{tab:res21} price impact\\
Rationing & Pro rata with a small random component & Removes the timing
channel, which is the one that would undo the mechanism; the random
component blunts the size-inflation incentive without imposing full
execution uncertainty & \Cref{tab:rationingres}\\
Settlement & Net settlement with per-fill audit records retained;
roll-over residual with standard time-in-force & Netting is a cost
saving with no incentive consequence provided the audit trail is kept
separately & Residual share, \Cref{tab:res23}\\
\bottomrule
\end{tabularx}
\end{table}

\subsection{How the Recommendation Changes by Venue}

The baseline is not universal, and three departures follow directly from
the analysis.

\textbf{Thin instruments.} The binding constraint is batch thickness,
not speed neutralisation. A longer interval, or a volume-triggered
close, keeps batches populated at the cost of longer delay. The
degeneracy diagnostic is the share of empty batches, which is why
\Cref{tab:sweep} reports it.

\textbf{Decentralised venues.} The interval is not a free parameter,
because settlement is already discrete and the natural boundary is the
block. Disclosure is also constrained, since submitted intents are
typically public. What the venue gains in exchange is the multi-asset
formulation of \eqref{eq:multiflow}: coherent cross-rates across all
pairs in a batch, which removes cross-pair arbitrage as a rent rather
than paying arbitrageurs to enforce consistency
\citep{walther2021multitoken,ramsay2023speedex}.

\textbf{Venues with a large retail segment.} The Glosten--Milgrom
argument for segregation is strongest here, because the uninformed
subset is both large and identifiable. The countervailing consideration
is regulatory rather than economic: participants in different segments
receive different prices for the same asset at the same instant, which
is difficult to reconcile with a uniform best-execution standard.

\subsection{The Unresolved Question: Maker and Taker}

\Cref{sec:segregation} argued that the maker--taker distinction has no
natural definition under simultaneous uniform-price clearing, and set
out four candidate redefinitions. This is not a peripheral issue. Fee
and rebate schedules are a substantial share of exchange revenue and a
substantial input into market maker economics, and a venue that adopts
batching without resolving the definition will find that its rebate
schedule is subsidising whichever behaviour its chosen proxy happens to
track.

Of the four candidates, persistence across batches has the strongest
economic claim, because it is the only one that measures the standing
provision of liquidity that the rebate is meant to compensate, rather
than an instantaneous property of a single order. It is also the only
one that cannot be satisfied by an order that exists for a single batch
and is designed to capture the rebate. The cost is that it requires the
engine to track order lineage across batches, which interacts with the
residual policy of \Cref{sec:settlement}: an order can only accrue
persistence if unmatched remainders roll over rather than being
cancelled.

\section{What This Design Can and Cannot Establish}
\label{sec:testable2}

\Cref{tab:predictions} classified nine predictions by testability, and
the distinction deserves restating in the discussion because it bounds
every claim in this thesis.

\textbf{Mechanical predictions are established.} P1 (zero
intra-interval dispersion), P4 (the delay lower bound), P8 (the effect
of $\tau$ on batch composition), and P9 (positive surplus relative to
submitted limits) follow from the clearing rule applied to a given order
set. The simulation confirms them, and where confirmation is by
construction rather than by measurement, \Cref{ch:results} says so.

\textbf{One prediction is partly reachable.} P2, the claim that batching
reduces the adverse selection component of the effective spread, is
partly testable. The mechanical part is: given this flow, does clearing
it in batches produce less price impact than clearing it continuously?
The behavioural part is not: liquidity providers in the replay cannot
requote in response to knowing they will not be sniped. A measured
reduction in price impact is therefore a lower bound on the effect
theory predicts, since the equilibrium response would go in the same
direction.

\textbf{Behavioural predictions are out of reach.} P3, P5, P6, and P7
all require participants who can adapt. Quoted spreads narrow because
market makers requote; size inflation occurs because participants
respond to pro rata; boundary concentration occurs because participants
respond to a predictable gate; bid shading occurs because participants
respond to price uncertainty. None of these responses can occur in a
replay of orders submitted under a different mechanism. What the
simulation provides instead is the mechanical baseline against which
such a response would have to be measured, which is a genuine
contribution but a different one.

The honest summary is that a replay design is the right instrument for
isolating what a mechanism does to a given order flow, and the wrong
instrument for establishing what a mechanism does to an equilibrium.
\Cref{sec:futurework} describes the extension that would address the
second question.

\section{Positioning Relative to the Literature}
\label{sec:positioning}

Against \citet{budish2015hft}, this thesis takes the mechanism as given
and asks how it is built. Their contribution is the equilibrium argument
for batching; the contribution here is the enumeration of what remains
undetermined once that argument is accepted, and the observation that
several of the remaining choices can undo a meaningful part of the
benefit. Price-time rationing inside the batch is the clearest example:
a venue can implement frequent batch auctions exactly as specified and
still have a speed race, because the allocation rule reintroduces one.

Against \citet{aquilina2022quantifying}, the relationship is
complementary. They measure what the continuous mechanism costs using
field data; this thesis measures what an alternative mechanism does to
the same flow. Their estimate is the benchmark against which any
measured reduction in adverse selection should be scaled, with the
caveat that their setting is an equity market with a different
participant population and a different latency distribution from a
crypto derivatives venue.

Against \citet{lee2026taiwan}, the relationship is methodological. Their
identification comes from a real transition, which means real behavioural
adjustment, and their external validity is correspondingly high; their
internal comparability is limited because everything moved at once. This
thesis has the opposite profile: perfect internal comparability, since
the flow is identical by construction, and limited external validity,
since the flow cannot respond. Neither design dominates, and read
together they bound the answer from two sides.

Against the decentralised exchange literature
\citep{walther2021multitoken,ramsay2023speedex,dutta2025cow}, the
contribution is to import the clearing formulation into a microstructure
framing. That literature solves the multi-asset clearing problem in
practice, including at speed, but evaluates the result mostly in terms
of throughput, gas cost, and solver incentives rather than in terms of
spreads and adverse selection. Placing those systems in the same design
space as an exchange matching engine, as \Cref{tab:casestudies} does,
makes it visible that they are instances of one mechanism rather than a
separate technology.

\section{Threats to Validity}
\label{sec:validity}

\paragraph{Internal validity.} The main threat is that a measured
difference reflects the implementation rather than the mechanism. This
is addressed by the shared type system and matching interface of
\Cref{lst:trait}, by computing all metrics offline from logs with one
shared implementation, by the invariant and unit-test layers of
\Cref{sec:validation}, and by reporting throughput and latency as
explicit controls. The degenerate-parameter check, in which the FBA
should converge to the CDA as $\tau \to 0$, is the strongest single
test, because it has a known answer and exercises the entire pipeline.

\paragraph{Construct validity.} Several metrics do not transfer
unchanged from a continuous to a batch setting, and each adaptation is a
place where the measured quantity may not be the intended one. The
quoted spread has no direct FBA analogue and is proxied by the marginal
unexecuted orders of the batch. The pre-trade midpoint referenced by
effective and realized spreads must be defined as the midpoint at the
start of the batch, which introduces an asymmetry: for the CDA the
reference is contemporaneous with the trade, for the FBA it is up to
$\tau$ old. Results should be read with this in mind, and the
sensitivity of the effective spread comparison to the choice of
reference point is worth reporting explicitly.

\paragraph{External validity.} Three limits apply. The data come from
one venue, and a crypto perpetuals exchange has a participant
population, a fee structure, and a latency profile that differ from an
equity market, so the magnitudes should not be transported. The sample
covers a limited number of sessions, which is why results are reported
per session as well as pooled. And, most fundamentally, the flow is
fixed: the mechanism cannot change behaviour, so any equilibrium effect
is absent by construction.

\paragraph{Data validity.} The dataset is captured at a single
non-validating node \citep{albers2026openbook}. What the node observed
is what the simulation replays, and any message the node missed or saw
in a different order than the venue's own engine is a discrepancy that
propagates into both arms of the comparison. Because it propagates into
both arms identically, it is largely differenced out of the paired
comparison, but it does affect the absolute levels and it affects the
CDA fidelity check of \Cref{sec:validation}.
